One of the consequences of the computable semantic model is that the meaning behind the execution of a software system is unambiguous. As a result, software systems can achieve successful coordination through shared meaning. I will first talk about how shared meaning leads to successful interaction in a general sense and then apply it to software systems.

Imagine you visit a library, find the book you are looking for and walk up to the librarian with the intention of checking out a book. In order to check out the book, you present the librarian with the book but the librarian was expecting a library card and the interaction halts. In this case, what it means to check out a book is different for the librarian and the visitor. Since there is no shared meaning between the two, the interaction fails and neither the librarian nor the visitor achieve their goal.

In this example, since it is two humans interacting, they can clarify and establish a shared meaning and move forward with the interaction. However, software systems can't improvise(yet), so shared meaning has to be established by construction. The CSM provides the mechanism to eliminate ambiguity in the meaning of interactions between software systems through their design.

\subsection{Semantically Compatible Designs}

Consider the system shown in Figure~\ref{fig:semantically_compatible_interactions}. This example establishes a distributed image compression service that connects multiple designs together. In this example, each node is its own design. When two designs are interacting with each other, in order to successfully coordinate they must share the same meaning.

\begin{figure}[htbp]
    \includesvg[
        width=0.99\linewidth
    ]{../assets/semantically_valid_interaction.svg}
    \caption{Semantically Compatible Interactions}
    \label{fig:semantically_compatible_interactions}
\end{figure}

When the meaning of the database is established, it also establishes the meaning of a database user. Therefore when the image compression service is interacting with the database, it takes

on the role of the database user and communicates with shared meaning. In this sense, the interaction between the two designs is semantically valid because they share the same meaning. In the diagram, I highlight this using a color scheme to indicate that each design is invoking shared meaning between both designs.

As a result, through semantically compatible interactions between designs, larger semantic worlds can be constructed. The larger semantic world establishes a new level of meaning that is defined by the semantically compatible designs that make up its reality. This frames a distributed system as a semantic world constructed from semantically compatible interactions between designs.

In order for successful interaction, the meaning that must be established on both sides of the interaction is unambiguous and this happens by design. In addition, the same principles that define the correctness of a design are equally as valid at this level of meaning. By tracking interactions between designs, environmental invariants can be identified at a particular level of meaning in the closed semantic world that represents a distributed system.

\subsection{Securing Interactions Between and Within Designs}

Finally, as a result, interaction between designs now has unambiguous meaning. There are meaningful implications for securing systems because the intention of the interaction is unambiguous. While before, the mechanical sequence had to be correct, now the design can ensure that it is the right narrative that is communicating with it. If the design is looked at as the brain of the software system, then this ensures that the correct brain is speaking, making it more difficult to compromise a system.

There are many approaches to implement this but one could be that the DAL can be extended to generate a dynamic checksum along the narrative path to identify unique narrative identity. There are more interesting ways that I can think of for securing it further by sharing the identity of the instance of the brain and using that to generate the proof. Regardless, the point here is, by eliminating ambiguity in a meaningful way, more automation will be enabled and securing the interaction between and within designs is another application for this framework.